%  =============================================================================
%  Datei-Hinweise & Konfiguration: siehe config.md
%  -> Abschnitt "anhang/quelltexte.tex"
%
%  Anhang "Projektstruktur und Quelltexte" (Leitfaden Kap. 5.9.2):
%  dokumentiert die Volltexte der Module des Handelsagenten (vgl. Verweis in
%  Kapitel 3, Abschnitt Implementierung) sowie die Projektstruktur als
%  Abbildung. Die Quelldateien liegen unverändert lauffähig unter
%  anhang/quelltexte/ und werden per \lstinputlisting eingebunden.
%
%  Druckaufbereitung (Transparenzhinweis, im Text ausgewiesen):
%    * Rufnummern des Messaging-Kanals sind redigiert (+49XXXXXXXXXXX).
%    * Typografische Sonderzeichen und Emoji-Symbole der Kommentar- und
%      Berichtstexte wurden fuer den Druck durch ASCII-Aequivalente ersetzt
%      (z. B. Pfeile -> , Warnsymbol [!]); die Programmlogik ist unveraendert.
%    * Umlaute/Eszett/Euro/Paragraf bleiben erhalten und werden ueber die
%      nachfolgende literate-Tabelle gesetzt.
%  =============================================================================

\chapter{Anhang: Projektstruktur und Quelltexte}
\label{cha:anhang-quelltexte}

% Literate-Tabelle: deutsche Sonderzeichen in Listings (pdflatex + listings
% verarbeiten UTF-8-Mehrbytezeichen nur ueber diese Zuordnung).
% Schriftgrad im Anhang: \footnotesize statt \small (ca. 7.000 Zeilen
% Quelltext; der Hauptteil behaelt die Vorgabe aus main.tex).
\lstset{%
    basicstyle=\ttfamily\footnotesize,
    literate={ä}{{\"a}}1 {ö}{{\"o}}1 {ü}{{\"u}}1
             {Ä}{{\"A}}1 {Ö}{{\"O}}1 {Ü}{{\"U}}1
             {ß}{{\ss{}}}1 {€}{{\euro{}}}1 {§}{{\S{}}}1
}

Dieser Anhang dokumentiert den Quelltextstand des in \cref{cha:umsetzung}
beschriebenen Handelsagenten zum 22.07.2026, unmittelbar vor der Aufnahme des
Echtgeldbetriebs. Abgedruckt sind die Volltexte sämtlicher Module des
Entscheidungs-, Order- und Betriebspfads sowie die zentralen
Konfigurationsdateien; damit ist jede in \cref{sec:umsetzung-implementierung}
beschriebene Absicherung im Quelltext nachvollziehbar. Die zwölf
Offline-Testsuiten und das bewusst vom Handelspfad entkoppelte öffentliche
Dashboard werden aus Platzgründen nicht im Volltext wiedergegeben, sondern in
der Projektstruktur ausgewiesen.

Für den Druck gelten drei redaktionelle Anpassungen, die die Programmlogik
unberührt lassen: Die Rufnummer des Messaging-Kanals ist redigiert
(\texttt{+49XXXXXXXXXXX}); Emoji-Symbole und typografische Sonderzeichen der
Kommentar- und Berichtstexte sind durch ASCII-Äquivalente ersetzt (etwa
\texttt{->} für Pfeile und \texttt{[!]} für das Warnsymbol); Zugangsdaten
erscheinen an keiner Stelle, da sie ausschließlich im verschlüsselten
Geheimnisspeicher liegen -- die Datei \texttt{agent.env} und der Inhalt von
\texttt{\textasciitilde/.fom-agent/} sind deshalb nicht Bestandteil des
Anhangs.

\section{Projektstruktur}
\label{sec:anhang-projektstruktur}

\Cref{fig:projektstruktur} zeigt die Struktur des Projekts. Es handelt sich um
einen einzigen Codestand, der auf beiden virtuellen Maschinen ausgerollt wird;
die Zuordnung der Module zu collector-vm und agent-vm ist je Datei vermerkt.
Fett gesetzte Dateien sind in den nachfolgenden Abschnitten im Volltext
abgedruckt.

\begin{figure}[htbp]
    \centering
    \includesvg[width=0.92\textwidth,height=0.88\textheight,keepaspectratio]{projektstruktur}
    \caption[Projektstruktur des Handelsagenten]{Projektstruktur des
        Handelsagenten: ein Codestand für beide virtuellen Maschinen
        (C = collector-vm, A = agent-vm).\quelleeigenki}
    \label{fig:projektstruktur}
\end{figure}

\section{Konfiguration}
\label{sec:anhang-konfiguration}

\lstinputlisting[language={},caption={\texttt{agentconfig.yaml} --
    Betriebskonfiguration: Strategien, Deployments, Sprachmodelle,
    Zeitbudgets, ML-Prüfinstanz},label={lst:agentconfig}]{anhang/quelltexte/agentconfig.yaml}

\lstinputlisting[language={},caption={\texttt{wknassign.yaml} -- zentrale
    Produktzuordnung Strategie zu Handelsinstrument},label={lst:wknassign}]{anhang/quelltexte/wknassign.yaml}

\section{Datenerfassung}
\label{sec:anhang-datenerfassung}

\lstinputlisting[language=Python,caption={\texttt{datacollect.py} --
    kontinuierliche Marktdatenerfassung mit Sanitätsfilter, Bar-Verdichtung
    und FRED-Crosscheck (collector-vm)},label={lst:datacollect}]{anhang/quelltexte/datacollect.py}

\lstinputlisting[language=Python,caption={\texttt{market\_calendar.py} --
    Börsenkalender: Handelstage, Feiertage,
    Sitzungsfenster},label={lst:marketcalendar}]{anhang/quelltexte/market_calendar.py}

\section{Strategieschicht}
\label{sec:anhang-strategieschicht}

\lstinputlisting[language=Python,caption={\texttt{strategy.py} --
    Strategie-Handler: Signalberechnung, DecisionRequest, Audit-Flags,
    Kostenkontext, Instrumentenkatalog},label={lst:strategy}]{anhang/quelltexte/strategy.py}

\lstinputlisting[language=Python,caption={\texttt{strategyloader.py} --
    Plugin-Lader: aktive Strategien, Deployments,
    ML-Konfiguration},label={lst:strategyloader}]{anhang/quelltexte/strategyloader.py}

\lstinputlisting[language=Python,caption={\texttt{plstrategy\_sma.py} --
    Plugin der SMA-Regimestrategie
    (Kreuzungssignal)},label={lst:plsma}]{anhang/quelltexte/plstrategy_sma.py}

\lstinputlisting[language=Python,caption={\texttt{plstrategy\_smacross.py} --
    Plugin-Variante des
    Kreuzungssignals},label={lst:plsmacross}]{anhang/quelltexte/plstrategy_smacross.py}

\lstinputlisting[language=Python,caption={\texttt{plstrategy\_smatrend.py} --
    Trendfolge-Alternative für den geplanten
    A/B-Vergleich},label={lst:plsmatrend}]{anhang/quelltexte/plstrategy_smatrend.py}

\lstinputlisting[language=Python,caption={\texttt{plstrategy\_buyhold.py} --
    Vergleichsmaßstab passives
    Halten},label={lst:plbuyhold}]{anhang/quelltexte/plstrategy_buyhold.py}

\lstinputlisting[language=Python,caption={\texttt{plstrategy\_momentum.py} --
    Vergleichsmaßstab
    Momentum},label={lst:plmomentum}]{anhang/quelltexte/plstrategy_momentum.py}

\section{Prüfinstanzen}
\label{sec:anhang-pruefinstanzen}

\lstinputlisting[language=Python,caption={\texttt{auditor.py} --
    Multi-LLM-Voting: Pre-Gate, unabhängige Voten, strikte Mehrheit,
    fail-safe Validator},label={lst:auditor}]{anhang/quelltexte/auditor.py}

\lstinputlisting[language=Python,caption={\texttt{mlforecast.py} --
    statistische Prüfinstanz: logistische Regression, Walk-Forward-Prüfung,
    Schatten- und Wirkmodus},label={lst:mlforecast}]{anhang/quelltexte/mlforecast.py}

\section{Orderausführung und Bankanbindung}
\label{sec:anhang-orderausfuehrung}

\lstinputlisting[language=Python,caption={\texttt{orchestrate.py} --
    abgesicherte Orderausführung: Frische-Gates, Intent-Journal,
    Orderbuch-Reconciliation,
    Circuit-Breaker},label={lst:orchestrate}]{anhang/quelltexte/orchestrate.py}

\lstinputlisting[language=Python,caption={\texttt{broker.py} --
    comdirect-Adapter: OAuth2, Session-TAN, Depot- und Orderbuchzugriff,
    tagesgültige Limit-Orders},label={lst:broker}]{anhang/quelltexte/broker.py}

\lstinputlisting[language=Python,caption={\texttt{credstore.py} --
    verschlüsselte Ablage der Zugangsdaten, getrennt je
    Konto},label={lst:credstore}]{anhang/quelltexte/credstore.py}

\lstinputlisting[language=Python,caption={\texttt{config.py} -- Lader der
    Umgebungskonfiguration},label={lst:config}]{anhang/quelltexte/config.py}

\section{Betrieb und Steuerung}
\label{sec:anhang-betrieb}

\lstinputlisting[language=bash,caption={\texttt{run\_cycle.sh} -- täglicher
    Entscheidungszyklus: Preflight, Strategie, Audit, Ausführung,
    Bericht},label={lst:runcycle}]{anhang/quelltexte/run_cycle.sh}

\lstinputlisting[language=Python,caption={\texttt{signal\_dispatcher.py} --
    Empfangsschleife des
    Messaging-Kanals},label={lst:dispatcher}]{anhang/quelltexte/signal_dispatcher.py}

\lstinputlisting[language=Python,caption={\texttt{control.py} --
    Befehlsprotokoll des Steuerkanals: Not-Aus, Broker-Login,
    ML-Umschaltung, Status},label={lst:control}]{anhang/quelltexte/control.py}

\lstinputlisting[language=Python,caption={\texttt{notify.py} -- Berichts- und
    Alarmversand mit Stale-Guard},label={lst:notify}]{anhang/quelltexte/notify.py}

\section{Auswertung}
\label{sec:anhang-auswertung}

\lstinputlisting[language=Python,caption={\texttt{evaluate\_votes.py} --
    wöchentliche Voting-Kalibrierung und ML-Auswertung gegen den
    Folgetags-Return},label={lst:evaluate}]{anhang/quelltexte/evaluate_votes.py}

%  =============================================================================
%  Ende
%  =============================================================================
